\subsection*{Procedimiento y operaciones del simulacro 2}
La geometría fija superficie y alargamiento antes de evaluar la polar. El equilibrio vertical fija el coeficiente de sustentación; de él salen incidencia, resistencia y potencia. Para buscar el mínimo se escribe $D=AV^2+B/V^2$, se deriva y se impone $AV^2=B/V^2$. Las raíces de equilibrio son las de la cuadrática en $V^2$ y exigen discriminante no negativo.

\begin{enumerate}

\item Área del trapecio.

\[\fitmath{S=\frac{12\cdot 2.1\cdot \left(1+0.5\right)}{2}\approx 18.9\,\mathrm{m^{2}}}\]

\item Alargamiento geométrico.

\[\fitmath{\mathrm{AR}=\frac{{12}^{2}}{S}}\]

\[\fitmath{\mathrm{AR}=\frac{{12}^{2}}{(18.9)}\approx 7.619048\,}\]

\item Coeficiente de resistencia inducida.

\[\fitmath{\mathrm{Ki}=\frac{1}{\pi\cdot 0.91\cdot \mathrm{AR}}}\]

\[\fitmath{\mathrm{Ki}=\frac{1}{\pi\cdot 0.91\cdot (7.619047619)}\approx 0.04591008\,}\]

\item Sumar término dependiente del perfil.

\[\fitmath{K=\mathrm{Ki}+0.007}\]

\[\fitmath{K=(0.04591007974)+0.007\approx 0.05291008\,}\]

\item Despejar eficiencia global; difiere de la eficiencia de sustentación.

\[\fitmath{\mathrm{eO}=\frac{1}{\pi\cdot \mathrm{AR}\cdot K}}\]

\[\fitmath{\mathrm{eO}=\frac{1}{\pi\cdot (7.619047619)\cdot (0.05291007974)}\approx 0.7896071\,}\]

\item ISA local.

\[\fitmath{T=288.15-0.0065\cdot 2500\approx 271.9\,\mathrm{K}}\]

\item Presión ISA.

\[\fitmath{p=101325\cdot {\left(\frac{T}{288.15}\right)}^{\frac{9.80665}{287.05287\cdot 0.0065}}}\]

\[\fitmath{p=101325\cdot {\left(\frac{(271.9)}{288.15}\right)}^{\frac{9.80665}{287.05287\cdot 0.0065}}\approx 74682.52\,\mathrm{Pa}}\]

\item Densidad local.

\[\fitmath{\rho=\frac{p}{287.05287\cdot T}}\]

\[\fitmath{\rho=\frac{(74682.51762)}{287.05287\cdot (271.9)}\approx 0.9568588\,\mathrm{kg/m^{3}}}\]

\item Escala de fuerza aerodinámica.

\[\fitmath{\mathrm{qS}=0.5\cdot \rho\cdot {72}^{2}\cdot S}\]

\[\fitmath{\mathrm{qS}=0.5\cdot (0.9568587862)\cdot {72}^{2}\cdot (18.9)\approx 46875.36\,\mathrm{N}}\]

\item Imponer sustentación igual al peso.

\[\fitmath{C_L=\frac{34000}{\mathrm{qS}}}\]

\[\fitmath{C_L=\frac{34000}{(46875.36371)}\approx 0.7253277\,}\]

\item Pendiente tridimensional del modelo lineal.

\[\fitmath{a=\frac{6}{1+\frac{6}{\pi\cdot 0.91\cdot \mathrm{AR}}}}\]

\[\fitmath{a=\frac{6}{1+\frac{6}{\pi\cdot 0.91\cdot (7.619047619)}}\approx 4.704183\,\mathrm{1/rad}}\]

\item Convertir el incremento angular a grados y sumar ángulo de sustentación nula.

\[\fitmath{\alpha=-2+\frac{\frac{C_L}{a}\cdot 180}{\pi}}\]

\[\fitmath{\alpha=-2+\frac{\frac{(0.7253277055)}{(4.704183392)}\cdot 180}{\pi}\approx 6.83431\,{}^{\circ}}\]

\item Parte inducida.

\[\fitmath{\mathrm{CDi}=\mathrm{Ki}\cdot {C_L}^{2}}\]

\[\fitmath{\mathrm{CDi}=(0.04591007974)\cdot {(0.7253277055)}^{2}\approx 0.02415331\,}\]

\item Parte de perfil dependiente de sustentación.

\[\fitmath{\mathrm{CDp}=0.007\cdot {C_L}^{2}}\]

\[\fitmath{\mathrm{CDp}=0.007\cdot {(0.7253277055)}^{2}\approx 0.003682702\,}\]

\item Coeficiente total.

\[\fitmath{C_D=0.026+\mathrm{CDi}+\mathrm{CDp}}\]

\[\fitmath{C_D=0.026+(0.02415330582)+(0.003682701962)\approx 0.05383601\,}\]

\item Equilibrio horizontal requiere este empuje.

\[\fitmath{D=\mathrm{qS}\cdot C_D}\]

\[\fitmath{D=(46875.36371)\cdot (0.05383600778)\approx 2523.582\,\mathrm{N}}\]

\item Potencia aerodinámica.

\[\fitmath{P=D\cdot 72}\]

\[\fitmath{P=(2523.582445)\cdot 72\approx 181697.9\,\mathrm{W}}\]

\item Comprobar Mach bajo.

\[\fitmath{\mathrm{Mach}=\frac{72}{\sqrt{1.4\cdot 287.05287\cdot T}}}\]

\[\fitmath{\mathrm{Mach}=\frac{72}{\sqrt{1.4\cdot 287.05287\cdot (271.9)}}\approx 0.2178126\,}\]

\item Margen de coeficiente respecto de pérdida.

\[\fitmath{\mathrm{CLmargin}=1.45-C_L}\]

\[\fitmath{\mathrm{CLmargin}=1.45-(0.7253277055)\approx 0.7246723\,}\]

\item Coeficiente del término cuadrático de velocidad.

\[\fitmath{A=0.5\cdot \rho\cdot S\cdot 0.026}\]

\[\fitmath{A=0.5\cdot (0.9568587862)\cdot (18.9)\cdot 0.026\approx 0.2351002\,\mathrm{kg/m}}\]

\item Coeficiente del término inverso cuadrático.

\[\fitmath{B=\frac{2\cdot K\cdot {34000}^{2}}{\rho\cdot S}}\]

\[\fitmath{B=\frac{2\cdot (0.05291007974)\cdot {34000}^{2}}{(0.9568587862)\cdot (18.9)}\approx 6764202\,\mathrm{kg\,m^{3}/s^{4}}}\]

\item Resolver la condición de derivada nula.

\[\fitmath{\mathrm{Vmin}={\left(\frac{B}{A}\right)}^{0.25}}\]

\[\fitmath{\mathrm{Vmin}={\left(\frac{(6764202.374)}{(0.2351002038)}\right)}^{0.25}\approx 73.23875\,\mathrm{m/s}}\]

\item Sustituir en resistencia.

\[\fitmath{\mathrm{Dmin}=2\cdot \sqrt{A\cdot B}}\]

\[\fitmath{\mathrm{Dmin}=2\cdot \sqrt{(0.2351002038)\cdot (6764202.374)}\approx 2522.114\,\mathrm{N}}\]

\item Empuje tras la reducción prescrita.

\[\fitmath{\mathrm{Td}=0.8\cdot D}\]

\[\fitmath{\mathrm{Td}=0.8\cdot (2523.582445)\approx 2018.866\,\mathrm{N}}\]

\item Discriminante negativo: ningún equilibrio nivelado del modelo.

\[\fitmath{\mathrm{Delta}={\mathrm{Td}}^{2}-4\cdot A\cdot B}\]

\[\fitmath{\mathrm{Delta}={(2018.865956)}^{2}-4\cdot (0.2351002038)\cdot (6764202.374)\approx -2285242\,\mathrm{N^{2}}}\]

\item Empuje de la última cuestión.

\[\fitmath{\mathrm{Tplus}=1.2\cdot \mathrm{Dmin}}\]

\[\fitmath{\mathrm{Tplus}=1.2\cdot (2522.114475)\approx 3026.537\,\mathrm{N}}\]

\item Discriminante del caso que sí admite raíces.

\[\fitmath{\mathrm{Dp}={\mathrm{Tplus}}^{2}-4\cdot A\cdot B}\]

\[\fitmath{\mathrm{Dp}={(3026.53737)}^{2}-4\cdot (0.2351002038)\cdot (6764202.374)\approx 2798867\,\mathrm{N^{2}}}\]

\item Raíz inferior positiva.

\[\fitmath{\mathrm{Vlo}=\sqrt{\frac{\mathrm{Tplus}-\sqrt{\mathrm{Dp}}}{2\cdot A}}}\]

\[\fitmath{\mathrm{Vlo}=\sqrt{\frac{(3026.53737)-\sqrt{(2798867.028)}}{2\cdot (0.2351002038)}}\approx 53.65332\,\mathrm{m/s}}\]

\item Raíz superior positiva.

\[\fitmath{\mathrm{Vhi}=\sqrt{\frac{\mathrm{Tplus}+\sqrt{\mathrm{Dp}}}{2\cdot A}}}\]

\[\fitmath{\mathrm{Vhi}=\sqrt{\frac{(3026.53737)+\sqrt{(2798867.028)}}{2\cdot (0.2351002038)}}\approx 99.97357\,\mathrm{m/s}}\]

\item Aplicar límite de sustentación a la primera raíz.

\[\fitmath{\mathrm{CLlo}=\frac{34000}{0.5\cdot \rho\cdot S\cdot {\mathrm{Vlo}}^{2}}}\]

\[\fitmath{\mathrm{CLlo}=\frac{34000}{0.5\cdot (0.9568587862)\cdot (18.9)\cdot {(53.65331944)}^{2}}\approx 1.306189\,}\]

\item Repetir para la segunda raíz; el equilibrio no prueba estabilidad.

\[\fitmath{\mathrm{CLhi}=\frac{34000}{0.5\cdot \rho\cdot S\cdot {\mathrm{Vhi}}^{2}}}\]

\[\fitmath{\mathrm{CLhi}=\frac{34000}{0.5\cdot (0.9568587862)\cdot (18.9)\cdot {(99.97356919)}^{2}}\approx 0.3762087\,}\]

\end{enumerate}
